% Generated from core-elements.md by md2tex.py, then reviewed against SKILL.md.

\subsection{Headings and inline text}
\label{sec:headings_and_inline_text}

\textbf{Objective:} check every core mapping in one file.

\keyterm{Pulse Code Modulation (PCM)} turns samples into codes. This is \emph{not} optional, and the signal \verb|data_out| leaves the module.

A run-in line with \emph{italic} text and a hard line break inside the paragraph.

\subsubsection{Lists}
\label{sec:lists}

\begin{itemize}
    \item First item
    
    \begin{itemize}
        \item Nested item
        
        \begin{itemize}
            \item Third level
        \end{itemize}
    \end{itemize}
    \item Second item
\end{itemize}

\begin{enumerate}
    \item Step one
    \item Step two
\end{enumerate}

\phantomsection\label{sec:deepest_heading}\textbf{Deepest heading:}

Text under a level-four heading.

\phantomsection\label{sec:is_this_a_question}\textbf{Is this a question?}

A run-in label that ends with a question mark gets no colon.

\needspace{10\baselineskip}
\subsection{Tables}
\label{sec:tables}

\begin{tabularx}{\textwidth}{|c|x{3.7cm}|L|}
    \hline
    \rowcolor{gray!15}
    \cc{ID} &
    \yy{Module}{3.7cm} &
    \CC{Description} \\ \hline
\endhead

    M01 &
    Control Unit &
    Manages instruction sequencing, state transitions, and coordination between datapath components during execution. \\ \hline

    M02 &
    Concatenation / Replication &
    Builds wider vectors from smaller pieces. \\ \hline

\rowcolor{white}
\caption{Tables}
\label{tab:tables}
\end{tabularx}

\begin{tabularx}{\textwidth}{|c|L|x{2.4cm}|}
    \hline
    \rowcolor{gray!15}
    \cc{Quantity} &
    \CC{Formula} &
    \yy{Symbols}{2.4cm} \\ \hline
\endhead

    Noise power &
    \tallcell{$\overline{q^2} = \dfrac{\Delta v^2}{12}$} &
    $\Delta v$ = step size\newline $q$ = error \\ \hline

    SQNR &
    $\text{SQNR} = 6.02n + 1.76$ dB &
    $n$ = bits \\ \hline

\rowcolor{white}
\caption{Tables}
\label{tab:tables_2}
\end{tabularx}

\subsection{Equations}
\label{sec:equations}

\begin{equation}
    f = \frac{1}{T}
    \label{eq:equations_1}
\end{equation}

where:

\begin{itemize}
    \item $f$ is frequency in Hz
    \item $T$ is the period in seconds
\end{itemize}

\begin{equation}
    \overline{q^2} = \frac{1}{\Delta v}\int_{-\Delta v/2}^{\Delta v/2} q^2\,dq
    \label{eq:equations_2}
\end{equation}

\begin{equation}
    = \frac{\Delta v^2}{12}
    \label{eq:equations_3}
\end{equation}

\begin{eqnote}{
\begin{equation}
    \boxed{B_{PM} = 2(\Delta f + B)}
    \label{eq:equations_4}
\end{equation}
}{\emph{Write also:} for FM, $\Delta f = k_f m_p/2\pi$; PM bandwidth grows with message frequency.}
\end{eqnote}

\subsection{Code}
\label{sec:code}

\begin{lstlisting}[language=bash, caption={Listing}]
make synth TOP=kronos_core
\end{lstlisting}

\begin{lstlisting}[caption={Listing}]
project/
+-- src/
|   +-- top.v
+-- tb/
\end{lstlisting}

\begin{lstlisting}[language=json, caption={Listing}]
{"design": "picorv32", "clock": 10}
\end{lstlisting}

\needspace{10\baselineskip}
\subsection{Figures}
\label{sec:figures}

\begin{figure}[H]
    \centering
    \includegraphics[width=0.75\linewidth]{figures/pre_de_emphasis_chain.png}
    \caption{Pre-emphasis and de-emphasis chain}
    \label{fig:pre_de_emphasis_chain}
\end{figure}

\textbf{Purpose:} shows where each filter sits.

\begin{figure}[H]
    \centering
    % TODO: insert graphic
    % \includegraphics[width=0.8\linewidth]{figures/missing.png}
    \caption{A figure that does not exist yet}
    \label{fig:missing}
\end{figure}

\subsection{Unicode and links}
\label{sec:unicode_and_links}

Bandwidth \texorpdfstring{$\approx$}{≈} 2(\texorpdfstring{$\beta$}{β} + 1)B, with \texorpdfstring{$\beta$}{β} \texorpdfstring{$\ge$}{≥} 0 and \texorpdfstring{$\theta$}{θ} \texorpdfstring{$\le$}{≤} \texorpdfstring{$\pi$}{π}. Signal \texorpdfstring{$\Rightarrow$}{⇒} noise, \okmark{} done, \nomark{} wrong, \warnmark{} careful, {\emojifont ★} key. Values: 10\texorpdfstring{$^{-6}$}{⁻⁶} s, x\texorpdfstring{$_{1}$}{₁}, ½ and 25 °C.

See \hyperref[sec:tables]{the tables section} and the \href{3-derivation-sheet.pdf\#sec:d1_bandwidth_of_pm}{derivation sheet}.
